
\documentclass{article}
\usepackage{colortbl}
\usepackage{makecell}
\usepackage{multirow}
\usepackage{supertabular}

\begin{document}

\newcounter{utterance}

\centering \large Interaction Transcript for game `dond', experiment `semi\_en', episode 19 with Qwen3{-}32B{-}t0.0.
\vspace{24pt}

{ \footnotesize  \setcounter{utterance}{1}
\setlength{\tabcolsep}{0pt}
\begin{supertabular}{c@{$\;$}|p{.15\linewidth}@{}p{.15\linewidth}p{.15\linewidth}p{.15\linewidth}p{.15\linewidth}p{.15\linewidth}}
    \# & $\;$A & \multicolumn{4}{c}{Game Master} & $\;\:$B\\
    \hline

    \theutterance \stepcounter{utterance}  
    & & \multicolumn{4}{p{0.6\linewidth}}{
        \cellcolor[rgb]{0.9,0.9,0.9}{
            \makecell[{{p{\linewidth}}}]{
                \texttt{\tiny{[P1$\langle$GM]}}
                \texttt{You are playing a negotiation game in which you have to agree on how to divide a set of items among you and another player.} \\
\\ 
\texttt{Rules:} \\
\texttt{(a) You and the other player are given a set of items. Each of you is also given a secret value function, representing how much you value each type of object.} \\
\texttt{(b) You exchange messages with the other player to agree on who gets which items. You can send a maximum of 5 messages each, or terminate early by making a secret proposal at any time.} \\
\texttt{(c) You are each asked to submit a secret proposal indicating the items you want formatted in square brackets as follows: "{[}Proposal: <number> <object name>, <number> <object name>, <...>{]}"} \\
\texttt{(d) If your proposals are complementary, i.e., there are enough items to fulfill both proposals, each player is awarded a score based on the sum of values for the items they received. Otherwise, both of you get zero points.} \\
\texttt{(e) Your goal is to maximize the score you receive.} \\
\\ 
\texttt{Let us start.} \\
\\ 
\texttt{The set of available items is:} \\
\\ 
\texttt{2 threads, 3 bolts, 1 picture} \\
\\ 
\texttt{Your secret value function is:} \\
\\ 
\texttt{thread: 0, bolt: 1, picture: 7} \\
\\ 
\texttt{IMPORTANT: Your messages, unless it is the secret proposal, are directly transmitted to the other player, so do not include any response to the rules or text announcing your message. To make a secret proposal, use the indicated format. Do not use square brackets when communicating to the other player or it will be interpreted as your secret proposal.} \\
            }
        }
    }
    & \\ \\

    \theutterance \stepcounter{utterance}  
    & \multicolumn{4}{p{0.6\linewidth}}{
        \cellcolor[rgb]{0.9,0.9,0.9}{
            \makecell[{{p{\linewidth}}}]{
                \texttt{\tiny{[P1$\rangle$GM]}}
                \texttt{<think>} \\
\texttt{Okay, let's see. I need to negotiate the division of 2 threads, 3 bolts, and 1 picture with another player. My value function is thread: 0, bolt: 1, picture: 7. So, the picture is super valuable to me, bolts are worth 1 each, and threads are worthless. The other player's value function is secret, so I don't know what they want. } \\
\\ 
\texttt{First, I should figure out what the other player might value. Since the items are threads, bolts, and a picture, maybe they value the threads or the picture as well. But since my picture is worth 7 to me, maybe I should try to get that. But if the other player also wants the picture, we might have a conflict. Alternatively, if they don't value the picture as much, maybe I can take it and give them bolts or threads.} \\
\\ 
\texttt{The goal is to maximize my score. Since the picture is worth 7, which is way higher than the bolts (each bolt is 1), I should prioritize getting the picture. But I need to make sure that the other player's proposal doesn't conflict with mine. If I take the picture, I can get 7 points, and then maybe take some bolts if possible. But there are 3 bolts and 2 threads. If I take the picture, maybe the other player can take the bolts and threads. But I need to check if their value function allows that.} \\
\\ 
\texttt{Wait, the other player's value function is secret. So I don't know if they value bolts or threads. But since there are 3 bolts and 2 threads, maybe they might prefer bolts. But if I take the picture, and they take the bolts and threads, maybe their proposal would be something like {[}Proposal: 3 bolts, 2 threads{]}. But if I take the picture, then the remaining items are 2 threads and 3 bolts. So if they want all the bolts and threads, that's possible. But if they also want the picture, then our proposals would conflict, and we both get zero. So I need to make sure that my proposal doesn't conflict with theirs.} \\
\\ 
\texttt{Alternatively, maybe I can propose taking the picture and some bolts. Let's see. If I take the picture (1) and maybe 1 bolt, then the remaining items would be 2 bolts and 2 threads. The other player could take those. But I} \\
            }
        }
    }
    & & \\ \\

    \theutterance \stepcounter{utterance}  
    & & & \multicolumn{2}{p{0.3\linewidth}}{
        \cellcolor[rgb]{0.9,0.9,0.9}{
            \makecell[{{p{\linewidth}}}]{
                \texttt{\tiny{[GM$|$GM]}}
                \texttt{proposal} \\
            }
        }
    }
    & & \\ \\

    \theutterance \stepcounter{utterance}  
    & & & \multicolumn{2}{p{0.3\linewidth}}{
        \cellcolor[rgb]{0.9,0.9,0.9}{
            \makecell[{{p{\linewidth}}}]{
                \texttt{\tiny{[GM$|$GM]}}
                [2, 3, 0]
            }
        }
    }
    & & \\ \\

    \theutterance \stepcounter{utterance}  
    & & \multicolumn{4}{p{0.6\linewidth}}{
        \cellcolor[rgb]{0.9,0.9,0.9}{
            \makecell[{{p{\linewidth}}}]{
                \texttt{\tiny{[GM$\rangle$P2]}}
                \texttt{You are playing a negotiation game in which you have to agree on how to divide a set of items among you and another player.} \\
\\ 
\texttt{Rules:} \\
\texttt{(a) You and the other player are given a set of items. Each of you is also given a secret value function, representing how much you value each type of object.} \\
\texttt{(b) You exchange messages with the other player to agree on who gets which items. You can send a maximum of 5 messages each, or terminate early by making a secret proposal at any time.} \\
\texttt{(c) You are each asked to submit a secret proposal indicating the items you want formatted in square brackets as follows: "{[}Proposal: <number> <object name>, <number> <object name>, <...>{]}"} \\
\texttt{(d) If your proposals are complementary, i.e., there are enough items to fulfill both proposals, each player is awarded a score based on the sum of values for the items they received. Otherwise, both of you get zero points.} \\
\texttt{(e) Your goal is to maximize the score you receive.} \\
\\ 
\texttt{Let us start.} \\
\\ 
\texttt{The set of available items is:} \\
\\ 
\texttt{2 threads, 3 bolts, 1 picture} \\
\\ 
\texttt{Your secret value function is:} \\
\\ 
\texttt{thread: 3, bolt: 0, picture: 4} \\
\\ 
\texttt{IMPORTANT: Your messages, unless it is the secret proposal, are directly transmitted to the other player, so do not include any response to the rules or text announcing your message. To make a secret proposal, use the indicated format. Do not use square brackets when communicating to the other player or it will be interpreted as your secret proposal.} \\
\\ 
\\ 
\texttt{The other player has submitted a secret proposal. It is now your turn to submit a proposal of your own.} \\
            }
        }
    }
    & \\ \\

    \theutterance \stepcounter{utterance}  
    & & & \multicolumn{4}{p{0.6\linewidth}}{
        \cellcolor[rgb]{0.9,0.9,0.9}{
            \makecell[{{p{\linewidth}}}]{
                \texttt{\tiny{[GM$\langle$P2]}}
                \texttt{<think>} \\
\texttt{Okay, let's see. The other player has already submitted a secret proposal, and now I need to respond with my own. The goal is to maximize my score based on my value function. My values are thread: 3, bolt: 0, picture: 4. The available items are 2 threads, 3 bolts, 1 picture.} \\
\\ 
\texttt{First, I need to figure out what the other player's proposal might be. Since their proposal is secret, I don't know exactly what they asked for. But maybe I can infer based on the items available and my own needs. Let's think about possible scenarios.} \\
\\ 
\texttt{If the other player is also trying to maximize their own score, they might have different valuations. For example, if they value bolts or the picture highly. But since I don't know their value function, I have to make assumptions. However, since the game requires our proposals to be complementary, I need to make sure that my proposal doesn't conflict with theirs.} \\
\\ 
\texttt{Wait, the rules say that if the proposals are complementary, meaning there are enough items to fulfill both, then we get our scores. Otherwise, both get zero. So my proposal must not require more of any item than is available when combined with theirs.} \\
\\ 
\texttt{Since I don't know their proposal, I need to make a proposal that is as flexible as possible. My highest value items are the picture (4) and threads (3 each). Bolts are worth 0 to me, so I don't care about them. So ideally, I want the picture and as many threads as possible.} \\
\\ 
\texttt{But there are only 2 threads and 1 picture. Let's consider possible allocations. If I take the picture, that's 4 points. Then, if I take both threads, that's 2*3=6, total 10. But if the other player also wants the picture, then we can't both have it. So maybe they have a proposal that includes the picture. If they do, I can't take it. But if they don't, I can.} \\
\\ 
\texttt{Alternatively, maybe the other player is taking some bolts. Since bolts are worth 0 to me, I don't mind if they take them. But if they take bolts, that leaves more threads and the picture for me. But how do I know?} \\
\\ 
\texttt{Since I have to make a secret proposal now, I need to decide what to ask for. The best strategy might be to ask for} \\
            }
        }
    }
     \\ \\

    \theutterance \stepcounter{utterance}  
    & & & \multicolumn{2}{p{0.3\linewidth}}{
        \cellcolor[rgb]{0.9,0.9,0.9}{
            \makecell[{{p{\linewidth}}}]{
                \texttt{\tiny{[GM$|$GM]}}
                \texttt{continue} \\
            }
        }
    }
    & & \\ \\

    \theutterance \stepcounter{utterance}  
    & & & \multicolumn{2}{p{0.3\linewidth}}{
        \cellcolor[rgb]{0.9,0.9,0.9}{
            \makecell[{{p{\linewidth}}}]{
                \texttt{\tiny{[GM$|$GM]}}
                \texttt{player was instructed to make a proposal} \\
            }
        }
    }
    & & \\ \\

\end{supertabular}
}

\end{document}
